% ===============================================================
%  PART V -- Simulation (Chapters 17--18)
%  The simulator, its verification, the workloads and the experiment
%  configurations. Measurements are in Part VI; code excerpts are in
%  Appendix E (rule 7).
% ===============================================================

\part{Simulation}
\label{part:simulation}

\begin{center}
\adjustbox{max width=\linewidth}{%
\begin{forest}
  part map,
  [Part V
    [Ch.~17: Simulator Design
      [Why our own simulator] [Traffic layer] [Timing layer] [Verification]]
    [Ch.~18: Workloads {\&} Configurations
      [Pattern kernels] [Synthetic mixtures] [Experiments] [Metrics]]
  ]
\end{forest}}
\end{center}

\begin{keybox}[What Part V does]
Chapter~17 describes the two-layer simulator we wrote in Python to apply the
method of Chapter~15. The traffic layer counts every message and byte under MSI,
MESI and MOESI. The timing layer adds per-core clocks, a shared bus and a memory
port, using the latencies assumed earlier in the report. The chapter ends by
checking the simulator against every analytical number of Parts~II to~IV.
Chapter~18 describes the workloads, the experiments and the metrics reported in
Part~VI. Chapter~10 planned to use the \texttt{gem5} full-system
simulator~\cite{binkert2011gem5} and OpenMP programs; \cref{sec:sim-why} explains
why we built the timing model ourselves and used synthetic workloads instead, and
\cref{sec:wl-scope} and Part~VII state what that costs.
\end{keybox}

% ===============================================================
%  Chapter 17 -- Simulator Design
% ===============================================================
\chapter{Simulator Design: A Two-Layer Coherence Simulator}
\label{ch:simulator}

\begin{figure}[htbp]
\centering
\adjustbox{max width=\linewidth}{%
\begin{forest}
  chapter map,
  [Chapter 17
    [Why our own simulator [Control] [Understanding] [Speed]]
    [Traffic layer [Caches {\&} holders] [Protocol rules] [Byte accounting]]
    [Timing layer [Report latencies] [Bus {\&} memory port] [Event loop]]
    [Verification [Analytical numbers] [Coherence checks] [Part IV identities] [Timing checks]]
  ]
\end{forest}}
\caption[]{Structure of Chapter~17: why we built a new
simulator, how its two layers work, and how it was checked before any result was
taken from it.}
\label{fig:ch17-map}
\end{figure}

% ---------------------------------------------------------------
\section{Why Build Our Own Simulator}
\label{sec:sim-why}
\label{sec:sim-arch}
\sectionmap{[Choice [Method needs one changed variable] [Team must understand every line]
  [Fast enough to sweep] [What we give up]]}
\index{MESI}\index{MOESI}\index{gem5}\index{simulator!traffic layer}\index{simulator!timing layer}%

Coherence studies usually use an established simulator such as
\texttt{gem5}~\cite{binkert2011gem5}, which models pipelines, caches, interconnects
and DRAM cycle by cycle, runs real programs and comes with several coherence
protocols. Three needs of this report pointed the other way. The method of
Chapter~15 requires MESI and MOESI to differ only in the three rules of
\cref{tab:transition-diff}, while the protocols shipped with a large simulator are
separate implementations that differ in many more ways, the situation of
\cref{fig:method-idea}(a); rewriting one to match the other would have been a
project of its own. Every result in Part~VI also had to be something we could
explain from first principles, which a few hundred lines whose rules are the tables
of this report allow and a large external code base does not. And testing H1 to H4
needs hundreds of paired runs, which our simulator finishes in about a minute on a
two-core machine (\cref{tab:sim-choice}). The price is realism: no pipeline, no
DRAM model and no real programs. We therefore report trends and approximate sizes,
not predictions of any real machine's run time.

\begin{table}[htbp]
\centering
\caption{The two options considered for the simulation study.}
\label{tab:sim-choice}
\small
\begin{tabular}{@{}L{3.3cm} L{4.6cm} L{4.6cm}@{}}
\toprule
 & \texttt{gem5}~\cite{binkert2011gem5} & Our two-layer simulator \\
\midrule
What is modelled & pipelines, caches, interconnect, DRAM, OS & private caches,
  one bus, one memory port, per-core clocks \\
Workloads & real compiled programs & synthetic traces of the four sharing patterns \\
MESI vs MOESI & two separate implementations & one implementation; three rules swapped \\
Output & cycles, latencies, energy & bytes, writebacks, cycles, AMAT, utilisation \\
Time to learn & weeks & the tables of this report \\
Run time of all experiments & days & about one minute \\
Satisfies \cref{fig:attribution-tree}? & no (protocols differ widely) & yes \\
\bottomrule
\end{tabular}
\end{table}

\Cref{fig:sim-arch} shows the structure. A workload generator produces a trace of
accesses. The traffic layer replays each access against per-core caches, applies
the protocol rules, updates states, counts messages and bytes, and checks
coherence. The timing layer around it decides when each core issues its next
access and uses the traffic layer's account of what happened to charge latency and
occupy shared resources. The traffic experiments use the traffic layer alone and
the timing experiments use both.

\begin{figure}[htbp]
\centering
\begin{bigfig}[0.92]
\begin{tikzpicture}[font=\sffamily\scriptsize,
    bx/.style={draw, thick, rounded corners=3pt, minimum height=0.8cm, align=center, inner sep=4pt},
    ar/.style={-{Stealth}, thick}]
  \node[bx, fill=tealC!14, draw=tealC, minimum width=2.4cm] (gen) at (0,0) {workload\\generator\\(Ch.~18)};
  \node[bx, fill=coreGrey!10, draw=coreGrey, minimum width=2.0cm] (tr) at (3.0,0) {trace\\(core, block, R/W)};
  \draw[thick, mesiBlue, rounded corners=4pt, fill=mesiBlue!4] (4.6,-2.2) rectangle (10.6,2.2);
  \node[font=\sffamily\bfseries\scriptsize, text=mesiBlue] at (7.6,1.95) {traffic layer (\texttt{coherence\_sim.py})};
  \node[bx, fill=white, draw=mesiBlue, minimum width=2.2cm] (cache) at (6.0,0.85) {per-core caches\\LRU, 64-B blocks};
  \node[bx, fill=white, draw=mesiBlue, minimum width=2.2cm] (hold) at (6.0,-0.85) {holders map\\block $\to$ copies};
  \node[bx, fill=moesiOrange!12, draw=moesiOrange, minimum width=2.4cm] (rules) at (9.0,0.85) {protocol rules\\MSI / MESI / MOESI};
  \node[bx, fill=memGreen!12, draw=memGreen, minimum width=2.4cm] (chk) at (9.0,-0.85) {coherence checks\\states + versions};
  \draw[ar, mesiBlue] (cache) -- (rules); \draw[ar, mesiBlue] (hold) -- (rules);
  \draw[ar, memGreen] (rules) -- (chk);
  \draw[thick, plumC, rounded corners=4pt, fill=plumC!4] (4.6,-5.0) rectangle (10.6,-2.7);
  \node[font=\sffamily\bfseries\scriptsize, text=plumC] at (7.6,-2.95) {timing layer (\texttt{timing\_sim.py})};
  \node[bx, fill=white, draw=plumC, minimum width=1.7cm] (q) at (5.6,-4.05) {event queue\\core clocks};
  \node[bx, fill=white, draw=plumC, minimum width=1.4cm] (bus) at (7.6,-4.05) {bus};
  \node[bx, fill=white, draw=plumC, minimum width=1.5cm] (mp) at (9.5,-4.05) {memory port\\owner ports};
  \node[bx, fill=amberC!14, draw=amberC, minimum width=2.4cm] (out) at (13.0,0) {counters\\bytes, writebacks,\\DSR, $F$};
  \node[bx, fill=amberC!14, draw=amberC, minimum width=2.4cm] (out2) at (13.0,-3.9) {cycles, AMAT,\\utilisation};
  \draw[ar] (gen) -- (tr); \draw[ar] (tr) -- (4.6,0);
  \draw[ar] (10.6,0) -- (out);
  \draw[ar, plumC] (10.6,-3.9) -- (out2);
  \draw[ar, plumC, {Stealth}-{Stealth}] (7.6,-2.2) -- node[right, font=\sffamily\tiny, align=left]{access issued /\\event record returned} (7.6,-2.7);
\end{tikzpicture}
\end{bigfig}
\caption[Simulator architecture]{The two layers. The traffic layer is the protocol:
it decides states, messages and bytes. The timing layer decides when accesses
happen and how long they take, from the event record returned for each access.}
\label{fig:sim-arch}
\end{figure}

% ---------------------------------------------------------------
\section{The Traffic Layer}
\label{sec:sim-traffic}
\sectionmap{[Traffic layer [Data structures] [Rules as code]
  [Byte accounting] [Event record]]}
\index{MESI}\index{MOESI}\index{MSI}\index{Owned state}\index{Exclusive state}\index{forced writeback}\index{snooping protocol}\index{LRU replacement}\index{owner}\index{simulator!timing layer}%

Each core's cache is an ordered dictionary from block number to a small record of
state and version. The dictionary order is the LRU order, so a hit moves the block
to the end and the victim in a full cache is the first entry. A second dictionary,
the holders map, records which cores hold each block and points at the same
records, so a snoop (``does any other cache hold this block, and in what state?'')
is one lookup instead of a search of every cache. Two more maps hold each block's
latest version and the version last written back, for the value checks of \cref{sec:method-correct}, and counters record every quantity of \cref{tab:metrics}.

The three protocols share one class, and two flags, ``has \stE'' and ``has \stO'',
select MSI, MESI or MOESI. With both flags set the rules are exactly those of
\cref{tab:transition-diff}. The read handler (\cref{lst:readmiss} in
\cref{app:code}) holds the only rule that separates MESI from MOESI on a remote
read: when another cache holds the block in \stM\ or \stO, the code either charges
a control and a data message and relabels the holder \stO\ (MOESI), or also charges
a writeback, updates memory's version and relabels the holder \stS\ (MSI and MESI).
One rule was not fixed by earlier chapters: a write miss to a block that another
cache holds dirty. The simulator moves the block from the dirty owner to the
writer, invalidates every other copy and writes nothing to memory, since the writer
becomes the new dirty owner. It does this the same way in all three protocols, so
it cannot favour either side, and none of the four sharing patterns of
\cref{tab:traffic-patterns} uses it, so it does not affect the analytical numbers
the simulator is checked against.

Every action is priced with the constants of \cref{tab:counting-rules}, written
once at the top of the source file so the simulator and the report cannot disagree
(\cref{tab:sim-bytes}). At the end of a run each block still in \stM\ or \stO\ is
charged one writeback (the drain), so a protocol cannot look cheaper by leaving
writebacks unpaid when measurement stops. Before each access the layer clears a
four-field event record and fills it as the rules run: the kind of access (hit,
read miss, write miss or upgrade), the supplier (a core, or memory), whether a
forced writeback happened, and whether a dirty block was evicted. The timing layer
reads only this record and never needs the protocol rules. Adding the record
changed no counter; the traffic results were byte-for-byte identical before and
after.

\begin{table}[htbp]
\centering
\caption{Bytes charged by the traffic layer for each action (interconnect /
memory).}
\label{tab:sim-bytes}
\small
\begin{tabular}{@{}L{6.6cm} c c c@{}}
\toprule
Action & MSI & MESI & MOESI \\
\midrule
Read miss served by memory & 80 / 0 & 80 / 0 & 80 / 0 \\
Read miss, block in \stM\ elsewhere & 152 / 64 & 152 / 64 & 80 / 0 \\
Read miss, block in \stO\ elsewhere & --- & --- & 80 / 0 \\
Write miss (any supplier) & 80 / 0 & 80 / 0 & 80 / 0 \\
Upgrade (write hit in \stS, or \stO) & 8 / 0 & 8 / 0 & 8 / 0 \\
Write hit in \stE & --- & 0 / 0 & 0 / 0 \\
Eviction of \stM\ (or \stO) & 72 / 64 & 72 / 64 & 72 / 64 \\
Drain: block still dirty at the end & 72 / 64 & 72 / 64 & 72 / 64 \\
\bottomrule
\end{tabular}
\end{table}

% ---------------------------------------------------------------
\section{The Timing Layer}
\label{sec:sim-timing}
\label{sec:sim-example}
\sectionmap{[Timing layer [Latencies from the report] [Three kinds of resource]
  [Event loop] [Serial or overlapped write back] [Two readers at once]]}
\index{MESI}\index{MOESI}\index{forced writeback}\index{snooping protocol}\index{bus utilisation}\index{owner}\index{simulator!traffic layer}\index{simulator!timing layer}%

The timing layer uses the latencies the report has already assumed and adds as few
as possible (\cref{tab:sim-latency}). Each added parameter is varied in at least one
experiment of Chapter~18, so no conclusion rests on one guessed value. Chapter~8
defined $T_{\mathrm{coh}}$ as the time to get a block from another cache
``including any writeback the protocol performs on the way'', which fits MESI's
dirty read. MOESI makes no writeback on the same read, so its latency is
$T_{\mathrm{coh}}$ minus the writeback transfer time, 51 cycles on the assumed bus.
This is the smallest change to the report's numbers that lets the two protocols
differ in latency at all.

\begin{table}[htbp]
\centering
\caption{Parameters of the timing layer. The upper block repeats the report's own
assumptions; the lower blocks are introduced or derived here.}
\label{tab:sim-latency}
\small
\begin{tabular}{@{}L{1.9cm} r L{5.6cm} L{3.4cm}@{}}
\toprule
Symbol & Value & Meaning & Source \\
\midrule
$T_{L1}$ & 1 & cache hit (cycles) & \cref{eq:amat}, Sarangi Ex.~147 \\
$T_{\mathrm{mem}}$ & 100 & read served by memory & \cref{eq:amat}, \cref{fig:hierarchy} \\
$T_{\mathrm{coh}}$ & 60 & read served by a modified copy, including the writeback & \cref{eq:amat-coh}, \cref{fig:amat-bars} \\
--- & --- & one message on the bus at a time; requests queue & \cref{eq:util} \\
\midrule
bus width & 8\,B/cycle & a 72-B message takes 9 cycles & assumed; Ch.~18 E1 \\
arbitration & 1 & cycles per message & assumed \\
$T_{\mathrm{upg}}$ & 10 & upgrade / invalidation round trip & assumed \\
$M_{\mathrm{occ}}$ & 10 & memory port busy per 64-B access & assumed; swept in E5 \\
$O_{\mathrm{occ}}$ & 10 & owner's cache port busy per block supplied & assumed; swept in E4 \\
think & 2--2000 & mean cycles of computation between memory operations & swept in E1 \\
\midrule
$t_{\mathrm{wb}}$ & 9 & writeback transfer: $\lceil 72/8 \rceil$ & derived \\
$T_{\mathrm{c2c}}$ & 51 & cache-to-cache read without writeback: $T_{\mathrm{coh}} - t_{\mathrm{wb}}$ & derived; swept in E4 \\
\bottomrule
\end{tabular}
\end{table}

Three kinds of shared resource are modelled, each serving one request at a time in
arrival order (\cref{fig:sim-resources}). The bus is held by each message for
$\lceil\text{bytes}/8\rceil + 1$ cycles, 2 for a control message and 10 for a data
or writeback message, so a MESI dirty read holds it for $2+10+10=22$ cycles and the
same read under MOESI for 12. The memory port is held for $M_{\mathrm{occ}}$ cycles
by each 64-byte read or write, including every forced writeback. Each core's cache
port is held for $O_{\mathrm{occ}}$ cycles whenever that core supplies a block as an
owner. The bus is split-transaction: it is busy only while messages are on it, not
while memory fetches the data. A dirty eviction is posted, so it occupies the bus
and memory port without stalling the core.

\begin{figure}[htbp]
\centering
\begin{bigfig}[0.9]
\begin{tikzpicture}[font=\sffamily\scriptsize,
    core/.style={draw, thick, rounded corners=2pt, minimum width=1.5cm, minimum height=0.75cm, align=center},
    res/.style={draw, thick, rounded corners=2pt, minimum height=0.7cm, align=center},
    ar/.style={-{Stealth}, thick}]
  \foreach \i/\x in {0/0, 1/2.4, 2/4.8, 3/7.2} {
    \node[core, fill=tealC!12, draw=tealC] (c\i) at (\x,2.6) {core P\i\\own clock};
    \node[res, fill=plumC!10, draw=plumC, minimum width=1.5cm] (p\i) at (\x,1.3) {cache port\\$O_{\mathrm{occ}}$};
    \draw[thick, tealC] (c\i) -- (p\i);
    \draw[thick, coreGrey] (p\i.south) -- (\x,0.25);
  }
  \draw[line width=4pt, mesiBlue!60] (-0.9,0) -- (10.4,0);
  \node[font=\sffamily\bfseries\scriptsize, text=mesiBlue, anchor=west] at (-0.9,-0.4) {bus: one message at a time, $\lceil \mathrm{bytes}/8\rceil + 1$ cycles each};
  \node[res, fill=amberC!14, draw=amberC, minimum width=2.0cm] (mp) at (9.6,1.3) {memory port\\$M_{\mathrm{occ}}$ per access};
  \node[res, fill=amberC!25, draw=amberC, minimum width=2.0cm] (mem) at (9.6,2.6) {memory\\$T_{\mathrm{mem}} = 100$};
  \draw[thick, amberC] (mem) -- (mp); \draw[thick, coreGrey] (mp.south) -- (9.6,0.25);
  \node[draw=alertRed, dashed, rounded corners=2pt, fit=(mp), inner sep=3pt] {};
  \node[text=alertRed, font=\sffamily\tiny, align=center] at (9.6,0.65) {every forced write back\\uses it too};
\end{tikzpicture}
\end{bigfig}
\caption[Timing-layer resources]{Each core has its own clock and a cache port used
when it supplies data as an owner; all cores share one bus and one memory port. A
request waits while the resource it needs is busy.}
\label{fig:sim-resources}
\end{figure}

The layer is a discrete-event simulation. Each core's accesses form a stream, and
between accesses the core computes for a think time drawn from an exponential
distribution. The draws come from a seed that depends only on the core and the
configuration, so the MESI and MOESI runs see identical think times. A priority
queue holds each core's next ready time. The loop takes the earliest core, lets the
traffic layer process its access, and charges the access: a
hit costs $T_{L1}$; a miss waits for and occupies the bus, then waits for the memory
port (memory supplies) or the owner's port (a cache supplies) and completes after
$T_{\mathrm{mem}}$, $T_{\mathrm{coh}}$ or $T_{\mathrm{c2c}}$; an upgrade completes
after $T_{\mathrm{upg}}$. The core's next access is ready when this one completes
plus the next think time. \Cref{lst:timing} in \cref{app:code} is the code that
charges a miss.

Whether the reader of a modified block waits for the forced writeback depends on
the bus, not the protocol. On some buses memory captures the flushed block as it
passes to the reader, so one transfer serves both; on others the writeback is a
separate step. The report's $T_{\mathrm{coh}}$, which includes the writeback,
implies the second. The timing layer supports both (\cref{fig:sim-wbmode}). In the
serial mode, the default, the reader waits $T_{\mathrm{coh}}=60$ cycles. In the
overlapped mode it waits $T_{\mathrm{c2c}}=51$, and the writeback occupies the bus
and memory port without delaying it. Part~VI reports both.

\begin{figure}[htbp]
\centering
\begin{bigfig}[0.95]
\begin{tikzpicture}[font=\sffamily\scriptsize, x=0.13cm]
  \foreach \y/\l in {0/{MESI, serial}, -1.1/{MESI, overlapped}, -2.2/{MOESI}} {
    \node[anchor=east] at (-2,\y) {\l};
  }
  \fill[mesiBlue!30] (0,-0.25) rectangle (2,0.25);
  \fill[plumC!35] (2,-0.25) rectangle (51,0.25);
  \fill[alertRed!45] (51,-0.25) rectangle (60,0.25);
  \node at (26,0) {owner lookup and data};
  \node[text=white, font=\sffamily\bfseries\tiny] at (55.5,0) {WB};
  \draw[very thick] (60,-0.35) -- (60,0.35) node[above]{60};
  \fill[mesiBlue!30] (0,-1.35) rectangle (2,-0.85);
  \fill[plumC!35] (2,-1.35) rectangle (51,-0.85);
  \fill[alertRed!20] (51,-1.35) rectangle (60,-0.85);
  \node at (26,-1.1) {owner lookup and data};
  \node[font=\sffamily\tiny, text=alertRed] at (69,-1.1) {WB off the critical path};
  \draw[very thick] (51,-1.45) -- (51,-0.75) node[above]{51};
  \fill[mesiBlue!30] (0,-2.45) rectangle (2,-1.95);
  \fill[plumC!35] (2,-2.45) rectangle (51,-1.95);
  \node at (26,-2.2) {owner lookup and data};
  \node[font=\sffamily\tiny, text=memGreen] at (63,-2.2) {no write back};
  \draw[very thick] (51,-2.55) -- (51,-1.85) node[above]{51};
  \draw[-{Stealth}] (0,-2.95) -- (80,-2.95) node[right]{cycles after bus grant};
  \foreach \x in {0,20,40,60} { \draw (\x,-2.9) -- (\x,-3.0) node[below]{\x}; }
\end{tikzpicture}
\end{bigfig}
\caption[Serial and overlapped writeback]{Latency of a read served by a modified
copy on an idle system. In serial mode the MESI reader waits for the 9-cycle
writeback; in overlapped mode it does not, but the writeback still uses the bus and
memory port. MOESI makes no writeback.}
\label{fig:sim-wbmode}
\end{figure}

On an idle system the latencies differ by 9 cycles, but waiting changes this.
Suppose P0 holds blocks A and B in \stM, and at the same moment P1 reads A and P2
reads B. Both need the bus, so one waits for the other. \Cref{fig:sim-gantt} and
\cref{tab:sim-gantt} trace the two reads in serial mode; the simulator gives exactly
these values. The first reader gains the 9 cycles of the writeback, as in
\cref{eq:dL}. The second gains 19: its own 9 plus the 10 cycles it no longer waits
for the first reader's writeback to clear the bus. With more cores competing, every
waiting request gains the same way, which is why H2 expects the time gain to grow
with utilisation.

\begin{figure}[htbp]
\centering
\begin{bigfig}[0.95]
\begin{tikzpicture}[font=\sffamily\scriptsize, x=0.125cm, y=0.62cm]
  \def\row#1#2{\node[anchor=east] at (-2,#1) {#2};}
  \row{0}{MESI: bus}
  \row{-1}{MESI: P1 reads A}
  \row{-2}{MESI: P2 reads B}
  \row{-3.6}{MOESI: bus}
  \row{-4.6}{MOESI: P1 reads A}
  \row{-5.6}{MOESI: P2 reads B}
  \fill[tealC!45] (0,-0.3) rectangle (22,0.3);  \node at (11,0) {P1: 2+10+10};
  \fill[plumC!40] (22,-0.3) rectangle (44,0.3); \node at (33,0) {P2: 2+10+10};
  \fill[tealC!20] (0,-1.3) rectangle (61,-0.7); \node at (30,-1) {served by P0 (incl.\ write back)};
  \draw[very thick] (61,-1.38) -- (61,-0.62) node[right, xshift=1pt, yshift=-6pt]{61};
  \fill[pattern=north east lines, pattern color=alertRed!70] (0,-2.3) rectangle (22,-1.7);
  \node[text=alertRed, fill=white, inner sep=1pt] at (11,-2) {wait 22};
  \fill[plumC!18] (22,-2.3) rectangle (83,-1.7); \node at (52,-2) {served by P0 (incl.\ write back)};
  \draw[very thick] (83,-2.38) -- (83,-1.62) node[right, xshift=1pt, yshift=-6pt]{83};
  \fill[tealC!45] (0,-3.9) rectangle (12,-3.3);  \node at (6,-3.6) {P1};
  \fill[plumC!40] (12,-3.9) rectangle (24,-3.3); \node at (18,-3.6) {P2};
  \fill[tealC!20] (0,-4.9) rectangle (52,-4.3); \node at (26,-4.6) {served by owner P0};
  \draw[very thick] (52,-4.98) -- (52,-4.22) node[right, xshift=1pt, yshift=-6pt]{52};
  \fill[pattern=north east lines, pattern color=alertRed!70] (0,-5.9) rectangle (12,-5.3);
  \node[text=alertRed, fill=white, inner sep=1pt] at (6,-5.6) {12};
  \fill[plumC!18] (12,-5.9) rectangle (64,-5.3); \node at (38,-5.6) {served by owner P0};
  \draw[very thick] (64,-5.98) -- (64,-5.22) node[right, xshift=1pt, yshift=-6pt]{64};
  \draw[-{Stealth}] (0,-6.6) -- (92,-6.6) node[right]{cycles};
  \foreach \x in {0,20,40,60,80} { \draw (\x,-6.55) -- (\x,-6.65) node[below]{\x}; }
  \draw[dashed, coreGrey] (0,-6.6) -- (0,0.4);
\end{tikzpicture}
\end{bigfig}
\caption[Two simultaneous dirty reads]{Two simultaneous reads of blocks held in \stM\
by P0 (serial mode). Under MESI each read holds the bus for 22 cycles because of the
writeback, so P2 waits 22 cycles; under MOESI it waits 12. Latencies include
$T_{L1}=1$.}
\label{fig:sim-gantt}
\end{figure}

\begin{table}[htbp]
\centering
\caption{The example of \cref{fig:sim-gantt} in numbers (cycles).}
\label{tab:sim-gantt}
\small
\begin{tabular}{@{}l c c c c@{}}
\toprule
 & P1 latency & P2 latency & Mean latency & Bus busy for the two reads \\
\midrule
MESI (serial) & 61 & 83 & 72 & 44 \\
MOESI & 52 & 64 & 58 & 24 \\
\midrule
MOESI saving & 9 & 19 & 14 (19\,\%) & 20 (45\,\%) \\
\bottomrule
\end{tabular}
\end{table}

The layer has no processor pipeline (a core stalls fully on a miss), no out-of-order
execution or prefetching, no DRAM banks or row buffers, one level of private cache,
and one snooping bus instead of a network. Execution time is the finish time of the
slowest core. These simplifications are why Part~VI compares MESI with MOESI
relative to each other on identical inputs and makes no claim about absolute run
times.

% ---------------------------------------------------------------
\section{Verification}
\label{sec:sim-verify}
\sectionmap{[Verification [Every analytical number reproduced] [Coherence checked on random runs]
  [Part IV identities hold] [Timing checks] [Only then results]]}
\index{simulator!traffic layer}\index{simulator!timing layer}\index{verification}%

We verified the simulator in four automated ways before taking any result from it;
the commands in \cref{lst:repro} (\cref{app:code}) rerun all of them in seconds.
First, the traffic layer was run on the pattern kernels of Chapter~18 and its byte
counts compared with the numbers derived by hand in Parts~II and~III
(\cref{tab:verify-patterns}); all 22 checks match exactly. Steady-state values are
measured as the difference between runs of $k+1$ and $k$ units, which cancels the
cold start and the drain.

\begin{table}[htbp]
\centering
\caption{Verification against \cref{tab:traffic-patterns} ($C=1$) and against
\cref{eq:upd-mesi,eq:upd-moesi} (bytes per round with $C$ consumers). Each cell of
the upper block is simulated\,/\,analytical bytes.}
\label{tab:verify-patterns}
\label{tab:verify-bc}
\small
\begin{tabular}{@{}l c c c c@{}}
\toprule
Pattern (unit) & MSI & MESI & MOESI & Result \\
\midrule
Private (read, write, evict) & 160 / 160 & 152 / 152 & 152 / 152 & 3 of 3 match \\
Read-only (per reader) & 80 / 80 & 80 / 80 & 80 / 80 & 3 of 3 match \\
Producer--consumer (per round) & 160 / 160 & 160 / 160 & 88 / 88 & 3 of 3 match \\
Migratory (per hand-off) & 160 / 160 & 160 / 160 & 88 / 88 & 3 of 3 match \\
\midrule
\multicolumn{5}{@{}L{12.4cm}@{}}{Producer--consumer with $C = 1, 2, 4, 8, 15$
  consumers: MESI simulated 160, 240, 400, 720, 1280\,B, equal to $160+80(C-1)$;
  MOESI simulated 88, 168, 328, 648, 1208\,B, equal to $88+80(C-1)$. 10 of 10
  match.} \\
\bottomrule
\end{tabular}
\end{table}

Second, the kernels exercise few transitions, so the layer was also run with every
check of \cref{sec:method-correct} enabled on 144 random workloads: every
combination of 2, 4, 8 and 16 cores, cache capacities of 2, 4 and 16 blocks and
unlimited, three sharing mixes and all three protocols, with 20\,000 accesses each.
The tiny caches force constant evictions of owned blocks, the case most likely to
expose a missing writeback. No state or value invariant was violated.

Third, the propositions of Part~IV were derived by hand, so the simulator checks
them and they check the simulator (\cref{tab:verify-identities}). Each held exactly,
to floating-point rounding, on all 468 paired runs of the two traffic experiments.
Fourth, the timing layer was checked against latencies worked out by hand on an
idle system and on one case with deliberate queueing (\cref{tab:verify-timing}); all
eleven checks pass.

\begin{table}[htbp]
\centering
\caption{The results of Part~IV checked on every paired traffic run (468 pairs).}
\label{tab:verify-identities}
\small
\begin{tabular}{@{}L{5.3cm} L{4.6cm} l@{}}
\toprule
Statement & Check on each pair & Result \\
\midrule
\Cref{prop:residency}: same contents & equal read misses, write misses,
  upgrades and hit rates & held in all pairs \\
\Cref{eq:dB-writebacks}: only writebacks differ & $\Delta B = 72\,(W_{\mathrm{MESI}} - W_{\mathrm{MOESI}})$ & held in all pairs \\
Forced writebacks are dirty reads & $B_{\mathrm{fwb}} = 72\,n_{\mathrm{dr}}$ under MESI & held in all pairs \\
\Cref{prop:never-worse}: bounds & $0 \le \Delta B \le 72\,n_{\mathrm{dr}}$ & held in all pairs \\
\Cref{eq:identity}: saving identity & $|S_B - \sigma\cdot72\,\dsr/\bar b| < 10^{-15}$ & held in all pairs \\
\bottomrule
\end{tabular}
\end{table}

\begin{table}[htbp]
\centering
\caption{Verification of the timing layer (cycles; latencies include $T_{L1}=1$).}
\label{tab:verify-timing}
\small
\begin{tabular}{@{}L{7.0cm} r r l@{}}
\toprule
Check & Simulated & Expected & \\
\midrule
Read miss served by memory & 101 & 101 & pass \\
MESI: two memory misses then an upgrade (sum) & 213 & 213 & pass \\
MESI serial: write miss + dirty read (sum) & 162 & 162 & pass \\
MOESI: write miss + dirty read (sum) & 153 & 153 & pass \\
MESI overlapped: write miss + dirty read (sum) & 153 & 153 & pass \\
Two simultaneous misses: mean latency & 107 & 107 & pass \\
Two simultaneous misses: mean bus wait & 6 & 6 & pass \\
Derived $T_{\mathrm{c2c}} = T_{\mathrm{coh}} - 9$ & 51 & 51 & pass \\
Bus cycles of a control and a data message & 2, 10 & 2, 10 & pass \\
MESI bus-busy cycles: write miss + dirty read & 34 & 34 & pass \\
MOESI bus-busy cycles: write miss + dirty read & 24 & 24 & pass \\
\bottomrule
\end{tabular}
\end{table}

% ===============================================================
%  Chapter 18 -- Workloads, Configurations and Metrics
% ===============================================================
\chapter{Workloads, Configurations and Metrics}
\label{ch:workloads}

\begin{figure}[htbp]
\centering
\adjustbox{max width=\linewidth}{%
\begin{forest}
  chapter map,
  [Chapter 18
    [Scope [Planned vs used] [What is lost]]
    [Kernels {\&} mixtures [Four micro-traces] [Knob $d$] [Seeds]]
    [Experiments [Traffic: 2 sweeps] [Timing: E1 to E5]]
    [Metrics {\&} reproducibility [Traffic] [Time]]
  ]
\end{forest}}
\caption[]{Structure of Chapter~18: the workloads, the
experiments that use them, and how every reported number is defined.}
\label{fig:ch18-map}
\end{figure}

% ---------------------------------------------------------------
\section{From the Problem Scope to the Workloads}
\label{sec:wl-scope}
\sectionmap{[Scope [Chapter 10 planned OpenMP and directories]
  [Used: synthetic patterns on a bus] [Why] [What is lost]]}
\index{MESI}\index{MOESI}\index{directory protocol}\index{OpenMP}%

\Cref{tab:scope-change} compares the scope set in \cref{sec:prob-scope} with what
the study uses. The changes have one cause. The method of Chapter~15 needs a
simulator in which MESI and MOESI differ only in three rules, and we chose to build
it rather than adapt a large external one. Running real programs would have meant
turning them into traces with an instrumentation tool we did not have time to set
up and validate, so we generated the four sharing patterns directly, since
Chapter~9 showed they decide the forced-writeback cost.

\begin{table}[htbp]
\centering
\caption{Scope planned in Chapter 10 and scope used in this study.}
\label{tab:scope-change}
\small
\begin{tabular}{@{}L{2.6cm} L{4.2cm} L{5.8cm}@{}}
\toprule
Item & Planned (\cref{sec:prob-scope}) & Used, and consequence \\
\midrule
Workloads & OpenMP C programs (e.g.\ matrix multiplication, image filtering) &
  synthetic traces of the four sharing patterns; results describe the patterns,
  not named programs \\
Organisation & snooping bus and directory & snooping bus only; directory
  results remain analytical (\cref{tab:p2p}) \\
Simulator & trace-driven, checked with \texttt{gem5} & our two-layer simulator;
  no cycle-accurate cross-check \\
Congestion metric & network-on-chip congestion & bus utilisation and queueing
  delay on the snooping bus (\cref{fig:res-h2}); no network-on-chip model \\
Cores & 2, 4, 8, 16 & 2, 4, 8, 16 (unchanged) \\
Protocols & MSI family & MSI, MESI, MOESI (unchanged) \\
Consistency, false sharing & out of scope & out of scope (unchanged) \\
\bottomrule
\end{tabular}
\end{table}

% ---------------------------------------------------------------
\section{Pattern Kernels and Synthetic Mixtures}
\label{sec:wl-kernels}
\label{sec:wl-mixtures}
\sectionmap{[Workloads [Four kernels for verification] [Four address pools]
  [One knob d for dirty sharing] [Fixed seeds]]}
\index{dirty sharing}\index{Dirty Sharing Ratio}\index{sharing patterns!producer--consumer}\index{sharing patterns!migratory}\index{verification}%

Four short deterministic traces reproduce one unit of each sharing pattern of
\cref{tab:patterns} on one block (\cref{fig:kernels}). They are used for
verification, and their names are reused for the pools of the synthetic generator.

\begin{figure}[htbp]
\centering
\begin{bigfig}[0.92]
\begin{tikzpicture}[font=\sffamily\scriptsize,
    rd/.style={draw=tealC, fill=tealC!15, thick, rounded corners=1pt, minimum width=0.85cm, minimum height=0.45cm},
    wr/.style={draw=alertRed, fill=alertRed!15, thick, rounded corners=1pt, minimum width=0.85cm, minimum height=0.45cm}]
  \node[anchor=east, font=\sffamily\bfseries\scriptsize] at (-0.3,0) {Private};
  \node[rd] at (0.6,0) {P0 R}; \node[wr] at (1.6,0) {P0 W};
  \node[anchor=west, text=coreGrey] at (2.3,0) {then evicted (drain); one core only};
  \node[anchor=east, font=\sffamily\bfseries\scriptsize] at (-0.3,-0.8) {Read-only};
  \foreach \i in {0,...,7} { \node[rd] at (0.6+\i*1.0,-0.8) {P\i\ R}; }
  \node[anchor=west, text=coreGrey] at (8.3,-0.8) {never written};
  \node[anchor=east, font=\sffamily\bfseries\scriptsize] at (-0.3,-1.6) {Producer--consumer};
  \node[wr] at (0.6,-1.6) {P0 W}; \node[rd] at (1.6,-1.6) {P1 R}; \node[rd] at (2.6,-1.6) {P2 R};
  \node at (3.35,-1.6) {$\cdots$}; \node[rd] at (4.1,-1.6) {P$C$ R};
  \node[wr] at (5.3,-1.6) {P0 W}; \node[rd] at (6.3,-1.6) {P1 R}; \node at (7.05,-1.6) {$\cdots$};
  \node[anchor=west, text=coreGrey] at (7.6,-1.6) {repeated rounds};
  \node[anchor=east, font=\sffamily\bfseries\scriptsize] at (-0.3,-2.4) {Migratory};
  \foreach \i/\c in {0/0,1/1,2/2,3/3} {
    \node[rd] at (0.6+\i*2.0,-2.4) {P\c\ R}; \node[wr] at (1.6+\i*2.0,-2.4) {P\c\ W}; }
  \node[anchor=west, text=coreGrey] at (8.3,-2.4) {block moves on};
\end{tikzpicture}
\end{bigfig}
\caption[The four pattern kernels]{The four pattern kernels, each on one block (R:
read, W: write). Steady-state costs come from differencing runs of $k+1$ and $k$
units.}
\label{fig:kernels}
\end{figure}

Programs mix the patterns, so the experiments use traces that mix them in
controlled proportions. The generator splits the address space into four pools
(\cref{fig:addr-pools}) and builds a trace one operation at a time. It picks a pool
with probability $(1-d)/2$ for private, $(1-d)/2$ for read-only, $d/2$ for
producer--consumer and $d/2$ for migratory, so $d$ is the share of operations in
the dirty-sharing patterns. A private operation is a random core accessing one of
its own 96 blocks, writing with probability 0.3. A read-only operation is a random
core reading one of 24 shared blocks. A producer--consumer operation picks one of 24
blocks, each with a fixed producer (block number modulo $N$); with probability
$p_w$ the producer writes it, and otherwise a random other core reads it. A
migratory operation has a core different from the last one read and then write one
of 12 blocks. The knob $d$ sets the input, and the DSR is measured. Because DSR is
weighted by transactions (\cref{eq:dsr-mix}) it rises steeply at small $d$, so the
sweeps use values of $d$ concentrated near zero. Every random choice is seeded by
the configuration, so a trace can be regenerated exactly (\cref{tab:wl-params}).

\begin{figure}[htbp]
\centering
\begin{bigfig}[0.9]
\begin{tikzpicture}[font=\sffamily\scriptsize]
  \def\h{0.75}
  \fill[skyC!25] (0,0) rectangle (7.2,\h);   \node at (3.6,\h/2) {private: 96 blocks per core ($96N$)};
  \fill[moesiOrange!30] (7.2,0) rectangle (9.0,\h); \node[align=center] at (8.1,\h/2) {prod.--cons.\\24};
  \fill[alertRed!25] (9.0,0) rectangle (10.2,\h); \node[align=center] at (9.6,\h/2) {migr.\\12};
  \fill[memGreen!25] (10.2,0) rectangle (12.0,\h); \node[align=center] at (11.1,\h/2) {read-only\\24};
  \draw[thick, coreGrey] (0,0) rectangle (12.0,\h);
  \foreach \x in {7.2,9.0,10.2} { \draw[thick, coreGrey] (\x,0) -- (\x,\h); }
  \node[anchor=north] at (3.6,-0.1) {probability $(1-d)/2$};
  \node[anchor=north, align=center] at (8.6,-0.1) {$d/2$ \;\; $d/2$};
  \node[anchor=north] at (11.1,-0.1) {$(1-d)/2$};
  \draw[decorate, decoration={brace, amplitude=5pt, mirror}] (7.2,-0.6) -- (10.2,-0.6)
     node[midway, below=6pt, text=alertRed]{dirty-sharing pools};
\end{tikzpicture}
\end{bigfig}
\caption[Address pools of the generator]{The generator's address space and the
probability of each pool. Only the two middle pools can produce dirty reads.}
\label{fig:addr-pools}
\end{figure}

\begin{table}[htbp]
\centering
\caption{Parameters of the workload generator.}
\label{tab:wl-params}
\small
\begin{tabular}{@{}L{3.6cm} L{3.6cm} L{5.4cm}@{}}
\toprule
Parameter & Values & Role \\
\midrule
Cores $N$ & 2, 4, 8, 16 & one stream per core \\
Dirty-sharing share $d$ & 0 to 1 (11 values, most below 0.1) & share of operations in the producer--consumer and migratory pools \\
Producer write probability $p_w$ & 0.15, 0.3, 0.5 & writes per round of producer--consumer sharing \\
Private write probability & 0.3 & writes to private data \\
Pool sizes & 96 per core / 24 / 12 / 24 blocks & private / prod.--cons. / migratory / read-only \\
Trace length & 60\,000 (traffic), 30\,000 (timing) & accesses per run \\
Seeds & 1, 2, 3 & three independent traces per configuration \\
\bottomrule
\end{tabular}
\end{table}

% ---------------------------------------------------------------
\section{Experiments, Metrics and Reproducibility}
\label{sec:wl-experiments}
\label{sec:wl-metrics}
\label{sec:wl-repro}
\sectionmap{[Experiments [{T1, T2 traffic}] [E1 to E5 timing]
  [Metrics defined once] [All paired] [About one minute to rerun]]}
\index{simulator!traffic layer}\index{simulator!timing layer}%

\Cref{tab:experiments} lists every experiment reported in Part~VI. Each
configuration runs under each protocol on the same trace, as Chapter~15 requires,
giving 468 paired traffic runs and 288 paired timing runs; the last column names
the hypothesis each experiment tests. \Cref{tab:metrics} defines every number
Part~VI reports. Traffic metrics come from the traffic layer and time metrics from
the timing layer, and protocols are always compared per pair (\cref{eq:paired}). The
code, inputs and outputs are kept together, and every number in Part~VI can be
regenerated with the commands in \cref{lst:repro} on any machine with Python~3.9 or
later and \texttt{numpy}, \texttt{pandas} and \texttt{matplotlib}. Seeds are fixed,
so repeated runs give byte-identical result files; we checked this for both layers.

\begin{table}[htbp]
\centering
\caption{Experiments. ``Fixed'' values hold unless the row varies them. Caches are
1024 blocks (64\,KB) unless varied, large enough for the private pool.}
\label{tab:experiments}
\small
\begin{tabular}{@{}l L{3.0cm} L{4.6cm} c L{2.0cm}@{}}
\toprule
 & Varies & Fixed & Pairs & Tests \\
\midrule
T1 & $N$, $d$ (11), $p_w$ (3), seed (3) & 64\,KB caches; MSI, MESI, MOESI & 396 & H1, H3, H4 \\
T2 & cache 4--1024 blocks, $d \in\{0.2,0.5,0.8\}$, seed & $N=8$, $p_w = 0.3$ & 72 & H3, H4 ($\sigma$) \\
\midrule
E1 & think 2--2000 cycles (10), $N \in\{8,16\}$, writeback mode & $d=0.5$, $p_w=0.3$ & 120 & H2 \\
E2 & $N$ (4), think $\in\{40, 300\}$ & $d=0.5$, $p_w=0.3$ & 24 & H1, H2 \\
E3 & cache 4--1024 blocks, $d \in\{0.2,0.5,0.8\}$ & $N=8$, think 40 & 72 & H4 \\
E4 & $T_{\mathrm{coh}}$ 30--200, $O_{\mathrm{occ}}$ 5--40 & $N=8$, $d=0.5$, think 40 & 60 & H4 (Effect B) \\
E5 & $M_{\mathrm{occ}}$ 5--40 & $N=8$, $d=0.5$, think 40 & 12 & H2 (memory) \\
\bottomrule
\multicolumn{5}{@{}l}{\footnotesize T = traffic layer only (trace order); E = both layers. All rows use 3 seeds.}
\end{tabular}
\end{table}

\begin{table}[htbp]
\centering
\caption{Metrics reported in Part~VI.}
\label{tab:metrics}
\small
\begin{tabular}{@{}L{2.4cm} L{6.4cm} L{3.8cm}@{}}
\toprule
Metric & Definition & Defined in \\
\midrule
$B_{\mathrm{tot}}$ & interconnect bytes, including the drain & \cref{tab:sim-bytes} \\
$B_{\mathrm{fwb}}$ & bytes of forced writebacks (MESI, MSI) & \cref{sec:prob-formal} \\
$F$ & $B_{\mathrm{fwb}}/B_{\mathrm{tot}}$ under MESI & \cref{eq:fwb-fraction} \\
$\dsr$ & dirty reads / coherence transactions, under MESI & \cref{sec:dsr-def} \\
$\bar b$ & $B_{\mathrm{tot}}/n_{\mathrm{tx}}$ under MESI & \cref{eq:identity} \\
$\sigma$ & saved epochs / dirty reads & \cref{eq:sigma} \\
$S_B$ & $1 - B_{\mathrm{tot}}(\mathrm{MOESI})/B_{\mathrm{tot}}(\mathrm{MESI})$ & \cref{eq:paired} \\
Memory bytes & bytes written to memory & \cref{tab:sim-bytes} \\
\midrule
$T$ & finish time of the slowest core (cycles) & \cref{sec:sim-timing} \\
AMAT & mean cycles per memory operation, hits included & \cref{eq:amat-coh} \\
$U_{\mathrm{bus}}$, $U_{\mathrm{mem}}$ & busy cycles / $T$ for bus and memory port & \cref{eq:util} \\
Speed-up & $T(\mathrm{MESI})/T(\mathrm{MOESI})$ & \cref{eq:paired} \\
Time reduction & $1 - T(\mathrm{MOESI})/T(\mathrm{MESI})$ & --- \\
Dirty-read latency & mean latency of reads served by a modified or owned copy & \cref{eq:dL} \\
\bottomrule
\end{tabular}
\end{table}
